論理を辿って一歩ずつ考える能力は、人類の知的活動を支える基盤です。
LLMは膨大な知識をもとに自然な文章を生成できますが、それは本当に「論理的に正しい」推論なのでしょうか。

私たちは、LLMに論理推論を教えるための研究を行っています。
その出発点としてまず、記号論理学の知見などを活用しつつ、論理推論の特徴を整理しました。
論理推論では、例えば「``A''という事実と、``AならばB''という事実から、``B''という事実を導く」というような推論規則に従って、新しい事実を導きます。
このような推論規則には多様な種類がありますが、それらは少数の「基礎的な」規則を出発点として導くことができます。

この考えに基づき、私たちは \textbf{LLMが論理推論を学ぶための「論理推論問題」}を設計しました。
与えられた事実に対して推論規則を一歩ずつ適用し、最終的な答えに到達する問題です。
問題設計では、基礎的な推論規則の組み合わせで解けることや、AやBに対して多様な事実を用いることで推論規則の一般性を学べること、等を重視しました。

さらに、こうした問題を大量に作成するための \textbf{自動生成アルゴリズム}を開発しました。
推論規則を多段階に積み重ねて問題の骨格を作り、語彙や自然言語テンプレートによって具体的な事実を生成します。

生成した問題を用いてLLMを学習させたところ、論理推論タスクの性能が大きく向上しました。
加えて、数学・科学・コーディングなど、幅広い推論タスクでも性能向上が確認されました。
これは、論理推論がさまざまな推論能力の基礎にあるためだと考えています。

本研究は2022年頃から継続しています。
学習データだけでなく、LLMの論理推論能力を評価するためのベンチマークも開発しています。
「それらしい答え」ではなく、真に論理的に正しい推論を目指して、研究を進めていきます。

\cvresearchfigure{data/research/FLD/main.pdf}
